\section{Triggers}

\subsection{Triggers de Auditoria}

\subsubsection{Trigger de Log Automatico}

\begin{lstlisting}[caption=Funcao generica de auditoria]
CREATE OR REPLACE FUNCTION audit_trigger()
RETURNS TRIGGER AS $$
BEGIN
    IF TG_OP = 'DELETE' THEN
        INSERT INTO audit.audit_logs (tabela, operacao, usuario_id, dados_antigos)
        VALUES (TG_TABLE_NAME, TG_OP, OLD.usuario_id, row_to_json(OLD));
        RETURN OLD;
    ELSIF TG_OP = 'UPDATE' THEN
        INSERT INTO audit.audit_logs (tabela, operacao, usuario_id, dados_antigos, dados_novos)
        VALUES (TG_TABLE_NAME, TG_OP, NEW.usuario_id, row_to_json(OLD), row_to_json(NEW));
        RETURN NEW;
    ELSIF TG_OP = 'INSERT' THEN
        INSERT INTO audit.audit_logs (tabela, operacao, usuario_id, dados_novos)
        VALUES (TG_TABLE_NAME, TG_OP, NEW.usuario_id, row_to_json(NEW));
        RETURN NEW;
    END IF;
    RETURN NULL;
END;
$$ LANGUAGE plpgsql;
\end{lstlisting}

\subsubsection{Aplicacao dos Triggers}

\begin{lstlisting}[caption=Criacao dos triggers de auditoria]
-- Trigger para transacoes
CREATE TRIGGER trg_audit_transacoes
    AFTER INSERT OR UPDATE OR DELETE ON transacoes
    FOR EACH ROW EXECUTE FUNCTION audit_trigger();

-- Trigger para contas
CREATE TRIGGER trg_audit_contas
    AFTER INSERT OR UPDATE OR DELETE ON contas
    FOR EACH ROW EXECUTE FUNCTION audit_trigger();

-- Trigger para usuarios
CREATE TRIGGER trg_audit_usuarios
    AFTER INSERT OR UPDATE OR DELETE ON usuarios
    FOR EACH ROW EXECUTE FUNCTION audit_trigger();
\end{lstlisting}

\subsection{Triggers de Negocio}

\subsubsection{Atualizacao Automatica de Saldos}

\begin{lstlisting}[caption=Trigger para atualizar saldo automaticamente]
CREATE OR REPLACE FUNCTION atualizar_saldo_trigger()
RETURNS TRIGGER AS $$
BEGIN
    IF TG_OP = 'INSERT' OR TG_OP = 'UPDATE' THEN
        -- Atualizar saldo da conta origem
        PERFORM atualizar_saldo_conta(NEW.conta_id);
        
        -- Se for transferencia, atualizar conta destino
        IF NEW.tipo = 'transferencia' AND NEW.conta_destino_id IS NOT NULL THEN
            PERFORM atualizar_saldo_conta(NEW.conta_destino_id);
        END IF;
        
        RETURN NEW;
    ELSIF TG_OP = 'DELETE' THEN
        -- Atualizar saldo da conta origem
        PERFORM atualizar_saldo_conta(OLD.conta_id);
        
        -- Se for transferencia, atualizar conta destino
        IF OLD.tipo = 'transferencia' AND OLD.conta_destino_id IS NOT NULL THEN
            PERFORM atualizar_saldo_conta(OLD.conta_destino_id);
        END IF;
        
        RETURN OLD;
    END IF;
    RETURN NULL;
END;
$$ LANGUAGE plpgsql;

-- Aplicar trigger
CREATE TRIGGER trg_atualizar_saldo
    AFTER INSERT OR UPDATE OR DELETE ON transacoes
    FOR EACH ROW EXECUTE FUNCTION atualizar_saldo_trigger();
\end{lstlisting}

\subsubsection{Validacao de Integridade}

\begin{lstlisting}[caption=Trigger para validacoes de integridade]
CREATE OR REPLACE FUNCTION validar_integridade_trigger()
RETURNS TRIGGER AS $$
BEGIN
    -- Validar que usuario pertence ao grupo da transacao
    IF NEW.grupo_id IS NOT NULL THEN
        IF NOT EXISTS (
            SELECT 1 FROM usuario_grupos 
            WHERE usuario_id = NEW.usuario_id 
              AND grupo_id = NEW.grupo_id 
              AND ativo = true
        ) THEN
            RAISE EXCEPTION 'Usuario nao pertence ao grupo especificado';
        END IF;
    END IF;
    
    -- Validar que conta pertence ao usuario ou grupo
    IF NOT EXISTS (
        SELECT 1 FROM contas c
        WHERE c.id = NEW.conta_id
          AND (c.usuario_id = NEW.usuario_id OR 
               (c.compartilhada = true AND c.grupo_id = NEW.grupo_id))
    ) THEN
        RAISE EXCEPTION 'Usuario nao tem acesso a conta especificada';
    END IF;
    
    RETURN NEW;
END;
$$ LANGUAGE plpgsql;

-- Aplicar trigger
CREATE TRIGGER trg_validar_integridade
    BEFORE INSERT OR UPDATE ON transacoes
    FOR EACH ROW EXECUTE FUNCTION validar_integridade_trigger();
\end{lstlisting}

\subsection{Triggers de Performance}

\subsubsection{Atualizacao de Views Materializadas}

\begin{lstlisting}[caption=Trigger para refresh de views materializadas]
CREATE OR REPLACE FUNCTION refresh_materialized_views()
RETURNS TRIGGER AS $$
BEGIN
    -- Refresh da view de resumo financeiro (em background)
    PERFORM pg_notify('refresh_mv', 'resumo_financeiro_mensal');
    
    RETURN COALESCE(NEW, OLD);
END;
$$ LANGUAGE plpgsql;

-- Aplicar trigger
CREATE TRIGGER trg_refresh_mv_transacoes
    AFTER INSERT OR UPDATE OR DELETE ON transacoes
    FOR EACH STATEMENT EXECUTE FUNCTION refresh_materialized_views();
\end{lstlisting}
